\documentclass[11pt,a4paper]{article}
\usepackage{fontspec}
\usepackage[french]{babel}
\usepackage[margin=2.3cm]{geometry}
\usepackage{amsmath,amssymb}
\usepackage{graphicx}
\usepackage{float}
\usepackage{booktabs}
\usepackage[font=small,labelfont=bf]{caption}
\usepackage[hidelinks]{hyperref}
\usepackage{xcolor}
\usepackage{listings}
\lstset{language=Python, basicstyle=\ttfamily\scriptsize, breaklines=true, columns=fullflexible,
  keepspaces=true, numbers=left, numberstyle=\tiny\color{gray}, frame=single, rulecolor=\color{gray!50},
  keywordstyle=\bfseries\color{blue!60!black}, commentstyle=\itshape\color{green!40!black},
  stringstyle=\color{red!60!black}, showstringspaces=false,
  literate={Δ}{{$\Delta$}}1 {α}{{$\alpha$}}1}

\graphicspath{{figures/}}
\newcommand{\ft}[1]{\widehat{#1}}
\setlength{\parskip}{0.4em}

\title{Édition d'image de Poisson par transformée de Fourier\\[0.2em]
\large Implémentation de Morel, Petro et Sbert (2012) et application à l'imagerie médicale}
\author{Célian Baffalie\\ Projet IMA, encadré par Elsa Angelini}
\date{Novembre 2026}

\begin{document}
\maketitle

\section{Introduction}

L'édition de Poisson modifie une image en imposant son gradient plutôt que ses valeurs :
on choisit un champ de vecteurs $V$ (le gradient « voulu ») et on cherche l'image $u$ dont le gradient
en est le plus proche. C'est ce qui permet de coller un objet sans laisser de couture, ou d'éclaircir
une zone sans créer de bord visible. L'article de Morel, Petro et Sbert \cite{morel} propose de résoudre
ce problème sur toute l'image, par transformée de Fourier, en une seule division par coefficient.
L'intérêt est que le coût ne dépend plus de la forme de la zone modifiée, qui peut alors être choisie
par un algorithme (un seuil, par exemple) plutôt qu'à la main.

J'ai réimplémenté la méthode en Python (NumPy) et reproduit toutes les expériences de l'article :
le rehaussement de contraste (§4.1), le collage et la suppression d'objet (§4.2, éq.~16), le mélange de
gradients (éq.~17) et les vérifications de la section~5. Ce rapport présente la méthode
(section~\ref{sec:methode}), les phénomènes que j'ai observés sur des images simples
(section~\ref{sec:phenomenes}), une comparaison avec les figures de l'article, sur ses propres images
(section~\ref{sec:article}), puis une application à des images médicales, en 2D et en 3D
(section~\ref{sec:medical}).

\section{La méthode}\label{sec:methode}

\subsection{Équation de Poisson}

Soit $f$ l'image de départ, définie sur un rectangle $R$, et $V=(V_1,V_2)$ le champ imposé. On cherche
\[
  u = \arg\min_u \int_R \lvert \nabla u - V \rvert^2 .
\]
L'équation d'Euler--Lagrange de ce problème est l'équation de Poisson $\Delta u = \operatorname{div} V$
dans $R$, avec la condition de Neumann $\frac{\partial u}{\partial n} = V\cdot n$ au bord, qui est la
condition « naturelle » de la minimisation. L'article prend un champ global :
$V = v$ dans la zone modifiée $\Omega$ et $V = \nabla f$ ailleurs (éq.~5). Il n'y a donc pas de condition
de Dirichlet sur le bord de $\Omega$, contrairement à l'édition de Poisson classique : on résout sur tout $R$.

\subsection{Résolution par Fourier}

L'image $H\times W$ est prolongée par symétrie miroir en une image $2H\times 2W$. Cette image est
périodique sans saut aux raccords, et sa dérivée normale est nulle sur les bords de $R$ : la condition
de Neumann est obtenue gratuitement. On calcule alors sa transformée de Fourier discrète
$\ft u_{mn}$ (éq.~9), avec $J=2H$, $L=2W$, et les pulsations $w_m = 2\pi m/J$, $w_n = 2\pi n/L$.
Dériver revient à multiplier par $i w$ (éq.~13), et l'équation de Poisson devient, pour chaque
coefficient,
\[
  \ft u_{mn} = \frac{-i\,\big(w_m \ft{V_1}_{mn} + w_n \ft{V_2}_{mn}\big)}{w_m^2 + w_n^2}
  \qquad\text{(éq.~12)}.
\]
Il reste à revenir en espace par transformée inverse et à garder le quart $H\times W$ de départ.
Le coût est celui de trois FFT, en $O(N\log N)$ pour $N$ pixels, quelle que soit la forme de $\Omega$.

\subsection{Quatre points absents de l'article}

\begin{itemize}
  \item \textbf{Recentrage des fréquences.} La FFT range les indices $m = 0,\dots,J-1$, mais les indices
  $m > J/2$ représentent les fréquences négatives $m-J$. Sans les ramener entre $-J/2$ et $J/2$ avant
  de multiplier par $i w_m$, le gradient obtenu n'est pas réel. L'article suppose implicitement la
  fonction « à bande limitée » (avant l'éq.~9), ce qui revient à ce recentrage.
  \item \textbf{Signes.} Avec les conventions de l'éq.~9, la dérivée multiplie par $+i w$ et la solution
  porte un signe $-i$ ; les équations 12 et 13 de l'article ont toutes deux le signe opposé, et les deux
  erreurs se compensent.
  \item \textbf{Fréquence nulle.} En $(m,n)=(0,0)$, l'éq.~12 donne $0/0$ : la constante n'est pas fixée
  par le gradient. Je prends $\ft u_{00} = \ft f_{00}$, c'est-à-dire la même moyenne que $f$. Pour un
  collage, je recale ensuite $u$ sur $f$ hors du masque (on ajoute la moyenne de $f-u$ hors de $\Omega$).
  \item \textbf{Retour dans $[0,255]$.} Après un rehaussement de contraste, $u$ sort de l'intervalle
  $[0,255]$ (par exemple de $-87$ à $396$ sur l'image de la figure~3 de l'article). L'article ne dit pas
  comment il affiche ses résultats, mais ils sont nettement plus clairs que $u$ écrêté : il y a donc une
  étape finale non décrite. J'utilise un étirement min-max (section~\ref{sec:etirement}).
\end{itemize}

\subsection{Extension aux volumes}

Rien dans la méthode n'est propre à la dimension 2. Pour un volume $H\times W\times D$, on prolonge par
symétrie sur les trois axes (volume $2H\times 2W\times 2D$), on utilise une FFT en 3D, et
\[
  \ft u_{mnp} = \frac{-i\,\big(w_m \ft{V_1} + w_n \ft{V_2} + w_p \ft{V_3}\big)}{w_m^2 + w_n^2 + w_p^2}.
\]

\subsection{Vérification (section 5 de l'article)}

Comme on résout sur toute l'image, il faut vérifier qu'avec $V=\nabla f$ on retrouve bien $f$
(problème~18 de l'article). Comme dans l'article, je recale $u$ pour qu'il ait la même moyenne et le même
écart-type que $f$, puis je mesure l'erreur quadratique moyenne (MSE), le PSNR et l'écart maximal (MD).
Le tableau~\ref{tab:valid} donne les résultats sur des images de la section~\ref{sec:medical}.

\begin{table}[H]
\centering\small
\begin{tabular}{@{}llrrr@{}}
\toprule
Image (taille) & Gradient & MSE & PSNR (dB) & MD \\
\midrule
Radio du thorax ($801\times700$) & spectral (éq.~13) & $8{,}8\cdot10^{-28}$ & 319 & $1{,}9\cdot10^{-13}$ \\
                                 & différences finies & $8{,}0$ & 39 & 125 \\
Radio de la main ($994\times700$) & spectral & $9{,}9\cdot10^{-28}$ & 318 & $2{,}6\cdot10^{-13}$ \\
                                  & différences finies & $13{,}5$ & 37 & 101 \\
IRM du cerveau ($700\times700$) & spectral & $2{,}1\cdot10^{-28}$ & 325 & $8{,}5\cdot10^{-14}$ \\
                                & différences finies & $5{,}0$ & 41 & 25 \\
\midrule
Article, table 1 & & $10^{-5}$ à $6\cdot10^{-3}$ & 69 à 96 & $0{,}005$ à $0{,}11$\\
\bottomrule
\end{tabular}
\caption{Reconstruction de $f$ à partir de son propre gradient. Avec le gradient spectral, l'erreur est
celle de l'arithmétique de la machine. Chaque résolution prend entre 0,3 et 0,9\,s.}
\label{tab:valid}
\end{table}

Avec le gradient spectral, la reconstruction est exacte : le gradient de l'éq.~13 et le solveur de
l'éq.~12 sont exactement inverses l'un de l'autre. L'article trouve des erreurs plus grandes, mais
invisibles ; il n'en donne pas l'origine (probablement des détails d'implémentation).

J'ai ajouté la ligne « différences finies » pour illustrer le §3.2 de l'article, qui insiste sur les
dérivées exactes. Si l'on donne au même solveur le gradient $f[k+1]-f[k]$, l'erreur atteint 100 niveaux
de gris sur les contours. En Fourier, cette différence multiplie par
$e^{iw}-1 = 2i\sin(w/2)\,e^{iw/2}$, et le solveur divise par $iw$ : on obtient
$\frac{\sin(w/2)}{w/2}\,e^{iw/2}\,\ft f$, c'est-à-dire l'image décalée d'un demi-pixel et un peu
floutée. Le gradient et le solveur doivent donc être cohérents.

\section{Phénomènes observés sur des images simples}\label{sec:phenomenes}

Pour comprendre la méthode, j'ai utilisé des images dont on connaît tout : un carré clair (200) sur un
fond sombre (50), de taille $128\times128$ (fichier \texttt{tests/test\_images\_simples.py}). Je l'ai
modifié de deux façons, chacune testant une question différente :
\begin{itemize}
  \item un bord net ou flou : comment le gradient spectral se comporte-t-il sur un saut ?
  (section~\ref{sec:gibbs}) ;
  \item un bruit : la dérivée amplifie-t-elle le bruit ? (section~\ref{sec:bruit}).
\end{itemize}
Dans tous les cas, $V=\nabla f$ redonne $f$ à $10^{-13}$ près : aucune de ces images ne gêne la
reconstruction. Ce qui change, c'est le gradient spectral lui-même, donc le résultat dès qu'on modifie ce
gradient, comme dans le rehaussement de contraste (section~\ref{sec:contraste}).

\subsection{Bord net : le gradient spectral oscille (effet de Gibbs)}\label{sec:gibbs}

Le gradient spectral est la dérivée exacte, aux pixels, de la courbe trigonométrique
$f(x) = \frac1N\sum_m \ft f_m e^{i w_m x}$ qui passe par tous les pixels. Pour un bord net, cette courbe
ondule entre les pixels (effet de Gibbs), et sa dérivée aussi. Pour une marche de hauteur $\Delta$ entre
les pixels 0 et 1, on peut le calculer : en espace, le gradient spectral est une convolution par le noyau
$(-1)^j/j$, et le pixel situé à $k\ge1$ pixels du bord vaut
\[
  f'[k] = \Delta\sum_{j\ge k}\frac{(-1)^{j+1}}{j},
  \qquad\text{soit}\quad \Delta\ln 2,\ \ \Delta(\ln2-1),\ \ \Delta\big(\ln2-\tfrac12\big),\ \dots
\]
(et de même de l'autre côté du bord). C'est un reste de la série harmonique alternée : il change de signe
à chaque pixel et ne décroît qu'en $\Delta/(2k)$. Pour $\Delta=150$, la formule donne $104$, $-46$, $29$,
$-21$, exactement les valeurs mesurées sur la figure~\ref{fig:gibbs}a. La différence finie, elle, vaut 150
entre les deux pixels du bord et 0 ailleurs.

\begin{figure}[H]
\centering
\includegraphics[width=0.85\linewidth]{gibbs.png}
\caption{Dérivée horizontale du carré, le long de la ligne 64. (a) Bord net : le gradient spectral oscille
de part et d'autre du bord, la différence finie non ; les croix sont les valeurs données par la formule
$\Delta\sum_{j\ge k}(-1)^{j+1}/j$. (b) Bord flouté ($\sigma=1$ pixel) : les deux coïncident.}
\label{fig:gibbs}
\end{figure}

Dès que le bord est un peu flou (figure~\ref{fig:gibbs}b), le spectre de l'image est presque nul près de
la fréquence maximale $\pi$ et les oscillations disparaissent. Pour le mesurer, j'ai pris la plus grande
valeur du gradient à plus de 4 pixels du bord, là où il devrait être nul :

\begin{center}\small
\begin{tabular}{@{}lccc@{}}
\toprule
 & Bord net & Flou $\sigma=1$ & Flou $\sigma=3$ \\
\midrule
Gradient spectral & 17,4 & 0,15 & 0,01 \\
Différence finie & 0 & 0 & 0 \\
\bottomrule
\end{tabular}
\end{center}

Ces oscillations ne gênent pas la reconstruction, car les valeurs alternées se compensent exactement.
Elles deviennent visibles dès que l'on modifie le gradient près d'un bord net, par exemple en le
multipliant par $\alpha$ d'un seul côté (section~\ref{sec:contraste}).

\subsection{Bruit : la dérivée l'amplifie}\label{sec:bruit}

J'ai ajouté au carré un bruit gaussien d'écart-type $\sigma=4$. Pour un bruit blanc,
la dérivée spectrale a un écart-type $\sigma\pi/\sqrt3\approx 1{,}81\,\sigma = 7{,}2$ (toutes les
fréquences jusqu'à $\pi$ sont multipliées par $w$) ; je mesure 7,04. La différence finie donne
$\sigma\sqrt2 = 5{,}66$ ; je mesure 5,58. Le gradient spectral amplifie donc un peu plus le bruit, mais la
reconstruction reste exacte ($1{,}4\cdot10^{-13}$). Le rehaussement de contraste ($T=100$, le fond est
sous le seuil) multiplie ensuite le bruit du fond par $\alpha$ : il passe de 4,00 à 10,01, soit
$\times 2{,}50$. C'est l'argument de l'article pour choisir $\alpha = 2{,}5$, qui fait passer un bruit
d'écart-type 4 à 10, limite jugée tolérable.

\subsection{Rehaussement de contraste : pourquoi une zone sombre peut devenir noire}\label{sec:contraste}

La méthode du §4.1 multiplie par $\alpha$ le gradient des pixels plus sombres qu'un seuil $T$ (éq.~15).
Sur la figure~\ref{fig:contraste_simple}, l'écart entre les deux carrés passe bien de 8 à 20.
Mais le fond monte de 170 à 209 et le grand carré descend de 35 à $-27$ : à l'affichage, il devient noir.

\begin{figure}[H]
\centering
\includegraphics[width=\linewidth]{contraste_simple.png}
\caption{Éq.~15 ($T=50$, $\alpha=2{,}5$) sur un fond clair (170) contenant un grand carré sombre (35) et
un petit carré à peine plus clair (43). (b) Le contraste du petit carré est bien multiplié par $2{,}5$,
mais tout le carré sombre descend sous 0 et s'affiche en noir. (c) Après l'étirement min-max
(section~\ref{sec:etirement}), il redevient gris foncé.}
\label{fig:contraste_simple}
\end{figure}

L'explication est le bord du carré. Le gradient spectral répartit le saut $\Delta = 135$ entre le carré et
le fond de part et d'autre du bord, symétriquement (section~\ref{sec:gibbs}) : la moitié du saut,
$\Delta/2$, est portée par les pixels du côté sombre. Ces pixels sont sous le seuil et leur gradient est
multiplié par $\alpha$. Le champ imposé demande donc un saut supplémentaire
$E = (\alpha-1)\Delta/2 \approx 101$, que la solution aux moindres carrés répartit entre les deux zones
proportionnellement à leurs aires : si $p$ est la part de l'image sous le seuil ($p=0{,}39$ ici), le fond
monte de $pE\approx 39$ et la zone sombre descend de $(1-p)E\approx 62$. L'article ne mentionne pas ce
problème : il est faible quand les zones sombres couvrent une grande partie de l'image (photo à
contre-jour) et fort pour un objet sombre entouré de clair.

On pourrait croire que le problème vient seulement des oscillations de Gibbs autour d'un bord net. Ce
n'est pas le cas : j'ai comparé une image à moitié sombre (40, avec des rayures d'amplitude 10) et à
moitié claire (160), avec une marche nette puis avec la même marche étalée sur une dizaine de pixels
($T=100$). La zone sombre descend de 40 à $-5{,}1$ avec la marche nette, et encore à $-0{,}9$ avec la
marche douce ; les rayures passent de 10 à 27 et 25. Avec une marche douce, tous les pixels de la marche
situés sous $T$ voient leur gradient multiplié par $\alpha$ : le saut supplémentaire vaut $(\alpha-1)$
fois la partie de la marche qui est sous le seuil, quelle que soit sa largeur. L'assombrissement fait donc
partie de la méthode, il ne disparaît pas en floutant l'image.

\subsection{Mon étirement min-max pour éviter les zones noires}\label{sec:etirement}

Comme $\ft u_{00}=\ft f_{00}$, l'image $u$ a la même moyenne que $f$ : ce que l'éq.~15 ajoute en contraste
doit être pris quelque part, et le sombre passe sous 0 pendant que le clair dépasse 255. Écrêter
($\min(\max(u,0),255)$) rend noires toutes les zones sous 0, c'est-à-dire exactement celles qu'on voulait
éclaircir. Je ramène donc $u$ dans $[0,255]$ en deux temps : un décalage, pour que le minimum soit 0, puis
une homothétie, pour que le maximum soit 255 :
\[
  v = \frac{255\,(u - \min u)}{\max u - \min u}.
\]
C'est une fonction affine croissante de $u$ : elle ne change pas l'ordre des niveaux de gris, elle multiplie
tous les contrastes par le même facteur $255/(\max u - \min u)$, et aucun pixel n'est écrêté. Sur la
figure~\ref{fig:contraste_simple}c, $u$ va de $-56$ à $228$ (facteur 0,90) ; le grand carré revient vers
26 et le fond monte à 238, et le petit carré reste bien distinct. C'est cet étirement que j'applique à
tous les résultats de contraste de la suite.

Il a deux défauts.
\begin{itemize}
  \item \textbf{Il dépend de deux pixels seulement}, le plus sombre et le plus clair. Or ces deux pixels
  sont souvent des dépassements de Gibbs autour d'un bord net, amplifiés par $\alpha$. Sur l'image de la
  figure~3 de l'article, $u$ va de $-87$ à $396$, alors que l'essentiel de l'image est entre 0 et 300 :
  le facteur vaut 0,53 et tous les contrastes de l'image sont presque divisés par deux, ce qui rend le
  ciel terne (figure~\ref{fig:article_contraste}). Sur l'IRM de la section~\ref{sec:medical}, le facteur
  tombe à 0,30 et l'image devient grise. Un étirement entre des percentiles (1\,\% et 99\,\%) serait plus
  robuste, au prix d'un peu d'écrêtage ; une autre possibilité, que j'ai essayée, est de garder la
  moyenne de $f$ et de diviser $u$ par l'amplitude attendue $\alpha T + (255-T)$, qui ne dépend d'aucun
  pixel extrême.
  \item \textbf{Il fait apparaître des taches de couleur.} Pour une image couleur, j'applique l'éq.~15 à
  la luminance, puis je multiplie les trois canaux par le rapport $v/f$. Pour un pixel presque noir, $f$
  est proche de 0 et ce rapport devient très grand : le moindre écart entre canaux, dû à la compression
  JPEG, devient une tache de couleur vive (figure~\ref{fig:taches}b).
\end{itemize}

\begin{figure}[H]
\centering
\includegraphics[width=\linewidth]{article_taches.png}
\caption{Berge du lac de la figure~5 de l'article, après l'éq.~15 et l'étirement. (b) Couleur obtenue en
multipliant les canaux par $v/f$ : taches dans les zones presque noires.}
\label{fig:taches}
\end{figure}

\subsection{Collage : ce qui se passe hors du masque}

La figure~\ref{fig:section5} reproduit l'expérience de la figure 9 de l'article. Il n'y a pas d'effet de
Gibbs hors du masque : l'écart décroît régulièrement avec la distance, sans oscillation. En revanche,
l'extérieur est modifié : 59 niveaux de gris au bord (aux points où les extrema de la texture rapide
touchent le bord, ce que l'article remarque aussi), 12 à 20 pixels et encore 7 à 60 pixels. Les écarts
rapides s'amortissent en quelques pixels ; le reste lent vient de ce que la variation lente de $f$ sous
le disque est remplacée par celle de $g$, et ce décalage à grande échelle est réparti des deux côtés du
bord, puisqu'il n'y a pas de condition de Dirichlet sur $\partial\Omega$. C'est la limite que l'article
signale lui-même au §6 : si l'extérieur ne doit pas changer du tout, l'ancienne méthode est préférable.

\begin{figure}[H]
\centering
\includegraphics[width=\linewidth]{section5.png}
\caption{Expérience de la figure 9 de l'article : une texture rapide (période 5 à 6 pixels) collée dans
un disque de rayon 60 sur une texture lente. (d) Écart maximal $|u-f|$ hors du disque, en fonction de la
distance au bord.}
\label{fig:section5}
\end{figure}

\subsection{Mélange de gradients (éq. 17)}

L'éq.~17 garde, dans $\Omega$, le gradient de $f$ là où $|\nabla f| > |\nabla g|$ et celui de $g$
ailleurs. À l'intérieur des lettres et dans le fond du texte, $\nabla g = 0$ et ce sont les rayures qui
restent ; au bord des lettres, $\nabla g$ est grand et impose le contour (figure~\ref{fig:melange}d).

\begin{figure}[H]
\centering
\includegraphics[width=\linewidth]{melange.png}
\caption{Texte (a) collé sur des rayures (b). (c) Éq.~16 : le gradient du texte est imposé dans tout le
rectangle et les rayures disparaissent. (d) Éq.~17 : en chaque pixel, on garde le gradient le plus fort
des deux images.}
\label{fig:melange}
\end{figure}

\section{Comparaison avec les figures de l'article}\label{sec:article}

Pour vérifier que mon code fait bien ce que fait l'article, je l'ai appliqué à ses propres images. Je les
ai extraites du PDF (figures~3 à 9) : elles font environ $420\times300$ pixels et sont compressées, donc
les comparaisons chiffrées ne donnent qu'un ordre de grandeur. Le code de cette section est dans le
fichier \texttt{tests/test\_images\_article.py}.

\paragraph{Reconstruction (table 1).} Sur les images des figures~3, 4, 5 et 7, la reconstruction de $f$ à
partir de son gradient donne une MSE de $10^{-27}$ (PSNR de 310 à 320\,dB, écart maximal
$5\cdot10^{-13}$), contre $10^{-5}$ à $6\cdot10^{-3}$ (69 à 96\,dB) dans l'article. Comme pour les images
médicales du tableau~\ref{tab:valid}, mon solveur est exact à la précision de la machine.

\paragraph{Rehaussement de contraste (figures 3 à 5).} Le masque de la zone sombre ($f<50$, calculé sur la
luminance) coïncide avec celui publié par l'article sur 97,5 à 99,2\,\% des pixels : le seuil et la
définition de la zone sont donc les mêmes. Le résultat publié est en revanche bien plus clair que $u$
(figure~\ref{fig:article_contraste}) : l'article applique une étape finale qu'il ne décrit pas. En
cherchant la meilleure fonction affine $a\,u+b$ qui approche le résultat publié, je trouve une pente de
0,70 à 0,76, ce qui confirme qu'il s'agit d'un étirement. Niveau moyen de la zone sombre :

\begin{center}\small
\begin{tabular}{@{}lccccc@{}}
\toprule
 & Original & Article & $u$ (éq.~15) & $u$ étiré (min-max) & Facteur de l'étirement \\
\midrule
Figure 3 (jockey) & 11 & 27 & $-21$ & 35 & 0,53 \\
Figure 4 (cascade) & 18 & 53 & $-3$ & 42 & 0,76 \\
Figure 5 (lac) & 11 & 38 & $-21$ & 61 & 0,74 \\
\bottomrule
\end{tabular}
\end{center}

Sans étirement, la zone sombre descend sous 0 et devient noire, alors qu'elle devait être éclaircie.
Avec l'étirement min-max, elle est éclaircie comme dans l'article, un peu plus pour le lac, un peu moins
pour la cascade. Sur la figure~3, le facteur de 0,53 rend le ciel plus terne que dans l'article : c'est le
premier défaut décrit à la section~\ref{sec:etirement}.

\begin{figure}[H]
\centering
\includegraphics[width=\linewidth]{article_contraste.png}
\caption{Figures 3 à 5 de l'article ($T=50$, $\alpha=2{,}5$). De gauche à droite : image d'origine, résultat
publié, $u$ affiché en écrêtant et $u$ après l'étirement min-max. Images tirées de \cite{morel}.}
\label{fig:article_contraste}
\end{figure}

\paragraph{Collage et mélange (figures 6 à 9).} J'ai refait les collages en retrouvant à la main les
masques et les décalages (pour la figure~6, le décalage du bateau est trouvé par corrélation : 40 pixels
vers le bas). Les résultats sont visuellement identiques à ceux de l'article
(figure~\ref{fig:article_collage}). Pour la figure~6, l'éq.~16 est plus proche du résultat publié
(30,6\,dB) que la copie des pixels (28,4\,dB). Pour la figure~8, l'éq.~17 (20,7\,dB) est plus proche que
l'éq.~16 (18,3\,dB), ce qui confirme que l'article a bien utilisé le mélange. Pour la figure~7, mon
résultat modifie moins l'extérieur du masque que celui de l'article (39,9\,dB contre 23,8\,dB par rapport
à l'image d'origine), sans doute parce que mon masque est plus serré que le sien.

\begin{figure}[H]
\centering
\includegraphics[width=\linewidth]{article_collage.png}
\caption{En haut, résultats publiés dans l'article (figures 6 à 9, images tirées de \cite{morel}) ; en bas,
mon code sur les mêmes images : clonage du bateau, effacement du faon (source : la même image décalée de 90
pixels), texte sur bois par mélange de gradients, texture sur texture.}
\label{fig:article_collage}
\end{figure}

\section{Application à l'imagerie médicale}\label{sec:medical}

\subsection{Rehaussement des zones sombres d'une radio du thorax}

Sur une radio du thorax, le seuil fixe $T=50$ de l'article ne sélectionne presque rien (3\,\% des
pixels) : il a été choisi pour des photographies. Le seuil automatique (25\,\% de l'image) donne $T=101$
et sélectionne les poumons (figure~\ref{fig:thorax}b). Après l'éq.~15, en 1,5\,s, la trame vasculaire des
poumons est nettement plus visible (figure~\ref{fig:thorax}e). On retrouve le défaut de la
section~\ref{sec:phenomenes} : la moyenne des poumons descend de 75 à 46, et les zones les plus sombres
passent sous 0 (4\,\% des pixels). Après l'étirement min-max, les poumons ne sont plus noirs (moyenne
109), mais $u$ va de $-195$ à $372$ : le facteur vaut 0,45 et l'image entière paraît grise et moins
contrastée que l'original, alors que le contraste local des poumons a augmenté. Les valeurs extrêmes
viennent des bords nets (inscriptions, bord du film). En imagerie médicale, il serait plus naturel de
laisser le radiologue choisir la fenêtre d'affichage, comme dans les logiciels de radiologie.

\begin{figure}[H]
\centering
\includegraphics[width=\linewidth]{thorax_contraste.png}
\caption{Radio du thorax ($801\times700$). Le seuil est choisi automatiquement pour que la zone sombre
couvre 25\,\% de l'image, comme le propose l'article : $T=101$, ce qui sélectionne les poumons. Le
résultat (c, e) est ramené dans $[0,255]$ par l'étirement min-max.}
\label{fig:thorax}
\end{figure}

Sur d'autres images (fond d'œil, radio de la main), la zone sous le seuil est surtout le fond noir
autour de l'organe, et la méthode apporte peu.

\subsection{Insérer une lésion}

Créer des images pathologiques réalistes à partir d'images saines est utile, par exemple pour augmenter
les données d'apprentissage d'un réseau de détection. La figure~\ref{fig:nodules} colle dans un poumon
sain une zone prise sur la radio d'un autre patient. La copie des pixels laisse un disque bien visible ;
avec l'éq.~16, les niveaux de gris se raccordent et la texture nodulaire semble appartenir à l'image,
mais les côtes de l'image de destination disparaissent dans le disque. Avec l'éq.~17, les côtes, dont le
gradient est plus fort que celui de la texture collée, sont conservées.

\begin{figure}[H]
\centering
\includegraphics[width=\linewidth]{nodules.png}
\caption{Insertion d'une zone de poumon présentant des nodules (a, radio d'une sarcoïdose) dans le poumon
droit d'une radio normale. (b) Copie des pixels. (c) Éq.~16. (d) Éq.~17.}
\label{fig:nodules}
\end{figure}

\subsection{Effacer un marqueur}

Les radios portent souvent des inscriptions (marqueurs, annotations). Comme dans la figure 7 de
l'article, on remplace le gradient d'un rectangle autour du marqueur par celui d'une zone voisine de la
même image (figure~\ref{fig:effacement}). Le résultat est sans couture, et le reste de l'image n'est
presque pas modifié (moins d'un niveau de gris sous la ligne 100), car la zone source et la zone
effacée ont le même contenu à grande échelle.

\begin{figure}[H]
\centering
\includegraphics[width=0.85\linewidth]{effacement.png}
\caption{Effacement des marqueurs de latéralité, en collant dans chaque rectangle le gradient de la zone
située 70 pixels à droite (éq.~16 avec la même image comme source, figure 7 de l'article).}
\label{fig:effacement}
\end{figure}

\subsection{Volumes 3D}

Le solveur 3D reconstruit un volume à partir de son gradient avec la même précision qu'en 2D
($1{,}4\cdot10^{-13}$). J'ai appliqué l'éq.~15 à une IRM cérébrale 3D, réduite d'un facteur 2 pour tenir
en mémoire, de deux façons : sur le volume entier, et coupe axiale par coupe axiale avec le solveur 2D.
Coupe par coupe, rien ne relie une coupe à la suivante et chacune reçoit son propre décalage : des stries
horizontales apparaissent sur les vues sagittale et coronale, alors qu'en 3D il n'y en a pas. Si l'on
traite un volume, il faut donc le traiter en 3D.

Le résultat n'est en revanche pas meilleur que l'original. Sur une IRM T1, la zone sous le seuil est
surtout l'air et le liquide cérébrospinal, qu'il n'y a pas d'intérêt à rehausser, et le tissu cérébral,
au-dessus du seuil, n'est pas modifié. De plus, l'étirement min-max a ici un facteur de 0,30 et rend
l'image entière grise. Le rehaussement de contraste de l'article n'est donc pas adapté à ce type d'image.

\section{Discussion}

La méthode de l'article est exacte, simple et rapide : trois FFT, quelle que soit la forme de la zone
modifiée, ce qui permet de la définir par un seuil. Elle se généralise sans difficulté aux volumes.
Mes essais font aussi apparaître ses limites :
\begin{itemize}
  \item le gradient spectral oscille autour des bords nets ; ces oscillations sont invisibles tant que le
  champ n'est pas modifié,
  mais apparaissent dès qu'on le modifie près d'un bord ;
  \item résoudre sur toute l'image modifie un peu l'extérieur de la zone éditée, d'autant plus que les
  contenus basse fréquence de la source et de la destination diffèrent ;
  \item le rehaussement de contraste enfonce les zones sombres entourées de clair, même quand le bord est
  doux, et le seuil $T=50$ n'est pas adapté aux images médicales (le seuil automatique l'est mieux) ;
  \item le rehaussement amplifie le bruit autant que le signal ;
  \item l'article ne dit pas comment il ramène le résultat dans $[0,255]$. Mon étirement min-max évite
  les zones noires et donne des images proches de celles de l'article, mais il dépend des deux pixels
  extrêmes, souvent des dépassements de Gibbs, et peut rendre toute l'image terne.
\end{itemize}

Sur ses propres images, mon code donne les mêmes résultats que l'article pour le collage et le mélange,
et une reconstruction plus précise (310 à 320\,dB contre 69 à 96\,dB).

\section{Conclusion}

J'ai réimplémenté la méthode de Morel, Petro et Sbert et reproduit toutes ses expériences, sur mes images
et sur celles de l'article. La reconstruction est exacte à la précision de la machine. J'ai expliqué les
phénomènes observés (oscillations du gradient spectral, amplification du bruit, zones sombres qui deviennent noires, perturbation hors du masque), proposé un
étirement min-max qui évite les zones noires, en précisant ses défauts, et appliqué la
méthode à des images médicales : rehaussement d'une radio du thorax, insertion d'une lésion, effacement de
marqueurs, puis extension aux volumes 3D, où le traitement en 3D évite les stries d'un
traitement coupe par coupe, même si le rehaussement n'améliore pas l'IRM. Une suite naturelle serait de
rendre l'étirement robuste aux pixels extrêmes (percentiles) et de corriger à la source le défaut des
zones sombres, par exemple en travaillant sur le logarithme de l'image.

\begin{thebibliography}{9}
\bibitem{morel} J.-M. Morel, A.~B. Petro, C. Sbert, \emph{Fourier implementation of Poisson image
editing}, Pattern Recognition Letters 33 (2012) 342--348.
\end{thebibliography}

\appendix
\section{Organisation du code}

\begin{itemize}
  \item \texttt{methodes/} : \texttt{solver.py} (éq.~9 à 13), \texttt{solver3d.py} (la même chose en 3D),
  \texttt{contraste.py} (éq.~15, en 2D et en 3D, étirement min-max, passage à la couleur),
  \texttt{collage.py} (éq.~16 et 17), \texttt{utils.py}.
  \item \texttt{tests/} : un fichier par expérience (reconstruction, contraste, images simples : bords
  et bruit, copie dans la même image ou depuis une autre image, mélange, section 5, solveur 3D,
  images de l'article).
  \item \texttt{docs/rapport/generer\_figures.py} produit toutes les figures et tous les chiffres de ce
  rapport.
\end{itemize}
Le code complet est reproduit à la fin de ce document (annexe~\ref{sec:code}). La répartition du travail
entre l'agent conversationnel et moi est détaillée à l'annexe~\ref{sec:agents}.

\section{Sources des images}

Images de Wikimedia Commons : radio du thorax normale et radio de la main (M.~Häggström, CC0) ; radio d'une
sarcoïdose (M.~Häggström, CC0) ; IRM cérébrale « MRI Brain T1 Axial (9) » (511KeV, CC BY-SA 4.0) ; fond
d'œil (domaine public).L'IRM 3D est le volume \texttt{sub-01\_T1w} du jeu de données OpenNeuro ds000001 (Schonberg et al.,
2012, CC0), lu avec nibabel. Les images de la section~\ref{sec:article} sont extraites des figures 3 à 9
de l'article \cite{morel} (© Elsevier), uniquement pour la comparaison ; elles ne sont pas redistribuées
avec le code.

\section{Utilisation des agents conversationnels}\label{sec:agents}

J'ai utilisé Claude (Anthropic), dans Claude Code, tout au long du projet, de trois façons.

\paragraph{Comprendre l'article.} Je m'en suis d'abord servi pour bien comprendre le papier de recherche :
je lui ai fait expliquer les équations, la transformée de Fourier discrète, le prolongement par symétrie
et la condition de Neumann, puis les phénomènes que j'observais (oscillations du gradient spectral,
zones sombres qui deviennent noires), jusqu'à ce que ce soit clair pour moi.

\paragraph{Imaginer et coder les fonctions principales.} Il m'a aidé à imaginer les fonctions à partir
des formules de l'article (le prolongement miroir, le gradient de l'éq.~13, le solveur de l'éq.~12, le
contraste de l'éq.~15, le collage et le mélange des éq.~16 et 17), puis à les coder dans le détail.
J'ai compris et retravaillé moi-même ces fonctions principales, rassemblées dans le dossier
\texttt{methodes/} : je sais expliquer chaque ligne, je les ai testées et modifiées. L'étirement min-max
de la section~\ref{sec:etirement} est une idée à moi, que j'ai ensuite implémentée avec son aide.

\paragraph{Les tests et les figures.} Je m'en suis encore plus servi pour les fichiers du dossier
\texttt{tests/} et pour le script qui produit les figures du rapport : je définissais l'expérience à faire
et ce que je voulais mesurer, et l'agent écrivait une grande partie du code correspondant, que je lançais
et dont j'analysais les résultats. Il m'a aussi aidé à rédiger et à mettre en forme ce rapport.

L'agent s'est parfois trompé (il a par exemple présenté comme venant de l'article des points qui n'y sont
pas) ; j'ai donc vérifié ses affirmations dans l'article et par des tests numériques.

\clearpage
\section{Code}\label{sec:code}

\subsection{Méthodes (dossier \texttt{methodes/})}
\lstinputlisting[title=\texttt{methodes/solver.py}]{../../methodes/solver.py}
\lstinputlisting[title=\texttt{methodes/solver3d.py}]{../../methodes/solver3d.py}
\lstinputlisting[title=\texttt{methodes/utils.py}]{../../methodes/utils.py}
\lstinputlisting[title=\texttt{methodes/contraste.py}]{../../methodes/contraste.py}
\lstinputlisting[title=\texttt{methodes/collage.py}]{../../methodes/collage.py}

\subsection{Tests (dossier \texttt{tests/})}
\lstinputlisting[title=\texttt{tests/test\_reconstruction.py}]{../../tests/test_reconstruction.py}
\lstinputlisting[title=\texttt{tests/test\_validation\_section5.py}]{../../tests/test_validation_section5.py}
\lstinputlisting[title=\texttt{tests/test\_images\_simples.py}]{../../tests/test_images_simples.py}
\lstinputlisting[title=\texttt{tests/test\_contraste\_exemples.py}]{../../tests/test_contraste_exemples.py}
\lstinputlisting[title=\texttt{tests/test\_contraste\_seuils.py}]{../../tests/test_contraste_seuils.py}
\lstinputlisting[title=\texttt{tests/test\_copie\_meme\_image.py}]{../../tests/test_copie_meme_image.py}
\lstinputlisting[title=\texttt{tests/test\_copie\_autre\_image.py}]{../../tests/test_copie_autre_image.py}
\lstinputlisting[title=\texttt{tests/test\_melange\_gradients.py}]{../../tests/test_melange_gradients.py}
\lstinputlisting[title=\texttt{tests/test\_solveur\_3d.py}]{../../tests/test_solveur_3d.py}
\lstinputlisting[title=\texttt{tests/test\_images\_article.py}]{../../tests/test_images_article.py}

\subsection{Figures du rapport}
\lstinputlisting[title=\texttt{docs/rapport/generer\_figures.py}]{generer_figures.py}

\end{document}
